\documentclass{tietreport}

% Essential packages
\usepackage{graphicx}
\usepackage{float}
\usepackage{amsmath}
\usepackage{amssymb}
\usepackage{booktabs}
\usepackage{tabularx}
\usepackage{array}
\usepackage[hidelinks]{hyperref}
\hypersetup{
    hidelinks,
    colorlinks=false,
    pdfborder={0 0 0}
}
\usepackage{listings}
\usepackage{xcolor}

\definecolor{codegreen}{rgb}{0,0.6,0}
\definecolor{codegray}{rgb}{0.5,0.5,0.5}
\definecolor{codepurple}{rgb}{0.58,0,0.82}
\definecolor{backcolour}{rgb}{0.97,0.98,0.99}
\definecolor{codeblue}{rgb}{0.13,0.29,0.53}

\lstdefinelanguage{JS}{
  keywords={const,let,var,function,if,else,return,await,async,new,true,false,null,document,window,localStorage},
  sensitive=true,
  comment=[l]{//},
  morecomment=[s]{/*}{*/},
  morestring=[b]',
  morestring=[b]",
  morestring=[b]`
}

\lstdefinestyle{jsstyle}{
    backgroundcolor=\color{backcolour},
    commentstyle=\color{codegreen}\itshape,
    keywordstyle=\color{codeblue}\bfseries,
    numberstyle=\tiny\color{codegray},
    stringstyle=\color{codepurple},
    basicstyle=\ttfamily\footnotesize,
    breaklines=true,
    captionpos=b,
    keepspaces=true,
    numbers=left,
    numbersep=6pt,
    showspaces=false,
    showstringspaces=false,
    tabsize=2,
    frame=single,
    rulecolor=\color{black!20},
    xleftmargin=12pt,
    framexleftmargin=10pt,
    aboveskip=8pt,
    belowskip=6pt,
    language=JS
}
\lstset{style=jsstyle}

\setlength{\headheight}{15pt}

% Bibliography configuration
\usepackage[
  date=year,
  style=alphabetic,
  backend=bibtex,
  sorting=nty,
  maxnames=5,
  minnames=3,
  maxalphanames=4,
  minalphanames=3,
  backref=true,
  doi=false,
  isbn=false,
  url=false,
  eprint=false
]{biblatex}
\addbibresource{references.bib}

% ---------------------------------------------------------------
% Helpers
% ---------------------------------------------------------------
% Highlighted reminder for content that still has to be filled in.
% Search for "TODO" in this file before submission.
\newcommand{\todo}[1]{\textcolor{red}{\textbf{[TODO: #1]}}}

% \shot[width]{file-in-images/}{caption}{label}
% Includes images/<file> if present; otherwise draws a visible placeholder
% so the report still compiles before all screenshots are added.
\newcommand{\shot}[4][0.9]{%
  \begin{figure}[H]
    \centering
    \IfFileExists{images/#2}{%
      \includegraphics[width=#1\textwidth]{images/#2}%
    }{%
      \fbox{\parbox[c][5cm][c]{0.8\textwidth}{\centering
        \textcolor{red}{Screenshot placeholder: \texttt{images/\detokenize{#2}}}}}%
    }
    \caption{#3}
    \label{#4}
  \end{figure}%
}

% \diagram[width]{file}{caption}{label}  (diagrams/<file>)
\newcommand{\diagram}[4][\textwidth]{%
  \begin{figure}[H]
    \centering
    \includegraphics[width=#1,height=0.72\textheight,keepaspectratio]{diagrams/#2}
    \caption{#3}
    \label{#4}
  \end{figure}%
}

% One full page per use case specification
\newcommand{\ucpage}[2]{%
  \begin{figure}[H]
    \centering
    \includegraphics[width=\textwidth,height=0.82\textheight,keepaspectratio]{diagrams/use_case_templates/#1}
    \caption{Use Case: #2}
  \end{figure}
  \clearpage
}

%% ==================================================
%% PROJECT-SPECIFIC FIELDS -- CHARGEV
%% ==================================================

\TitlePageHeader{%
  UCS503P - Software Engineering Project%
}

\ProjectTitle{%
  ChargEV: Find. Charge. Go Green.%
}

\ProjectSubTitle{%
  The Smart EV Charging Station Platform%
}

\author{%
  \texttt{1024170103} Pehal \\
  \texttt{1024170431} Briti Arora%
}

\TitlePageSubText{%
  Group: BE Third Year, CSE%
}

\AdvisorName{%
  Dr. Nisha Thakur%
}

\TitlePageFooterText{{%
    \large End Semester Evaluation \\[0.2cm]
    \bfseries Computer Science and Engineering
    Department%
  } \\[0.15cm]
  Thapar Institute of Engineering and Technology,
  Patiala%
}

%% ==================================================

\begin{document}

\maketitle

\tableofcontents

% =====================================================================
% CHAPTER 1: INTRODUCTION
% =====================================================================

\chapter{Introduction}

\section{Background and Context}

Electric vehicle (EV) adoption is accelerating, but the
charging infrastructure that supports it has not kept
pace in terms of usability. Drivers routinely do not
know which nearby stations are actually free, whether a
connector will be compatible with their vehicle, or
whether they will have to wait once they arrive. This
project, \textbf{ChargEV}, aims to remove this
uncertainty by giving EV drivers a single, reliable web
platform to discover charging stations, check real-time
availability, and reserve a slot in advance -- turning
charging from a source of anxiety into a predictable,
routine task.

The platform follows a simple user journey:
\textbf{Discover $\rightarrow$ Book $\rightarrow$ Charge
$\rightarrow$ Engage}. ChargEV combines location
services, real-time availability, online booking and a
gamified rewards layer (``EcoCoins'') in one interface.

\subsection{Problem Statement}

While EV charging infrastructure continues to expand,
the tools available to EV drivers to actually use that
infrastructure remain fragmented and unreliable. Current
challenges include:

\begin{itemize}
  \item \textbf{Discovery Gap:} It is hard to find nearby
    charging stations from a single, reliable source.
  \item \textbf{Unclear Availability:} Station status
    (free, occupied, out of service) is rarely known
    before arriving, so drivers waste time at busy
    stations.
  \item \textbf{No Advance Reservation:} Drivers have no
    convenient way to reserve a charging slot, and
    connector compatibility is often discovered only on
    arrival.
  \item \textbf{Generic Search Results:} Existing tools do
    not let drivers filter by connector type, distance,
    availability and price together.
  \item \textbf{Difficult Trip Planning:} Long-distance EV
    travel requires manually chaining together charging
    stops.
  \item \textbf{Low Engagement:} There is little
    incentive for drivers to keep returning to a single
    platform.
\end{itemize}

\section{Project Objectives}

To address these challenges, we have defined the
following core SMART (Specific, Measurable, Achievable,
Relevant, Time-bound) objectives:

\begin{enumerate}
  \item \textbf{Build a Centralized Station Locator:}
    Provide a single source of truth for nearby EV
    charging stations, searchable by location and
    filterable by connector type, distance,
    availability and price.
  \item \textbf{Deliver Real-Time Availability and Smart
    Booking:} Show live port availability and let
    drivers reserve a slot online so that a charge is
    guaranteed before they leave.
  \item \textbf{Provide Trip Planning and Navigation:}
    Use the Google Maps Platform to geocode locations,
    fetch place data and compute routes with charging
    stops.
  \item \textbf{Implement an Engagement Layer:} Reward
    sustainable usage with EcoCoins that can be
    redeemed as a discount at checkout, together with
    eco-impact tracking and badges.
  \item \textbf{Ensure Software Quality and Test
    Validation:} Specify all use cases formally,
    document the design with UML and data-flow models,
    and execute documented test cases against the
    implemented system.
\end{enumerate}

% =====================================================================
% CHAPTER 2: METHODOLOGY AND DESIGN
% =====================================================================

\chapter{Methodology and Design}

\section{System Architecture and Design Methodology}

ChargEV is designed as a multi-tier web application that
separates the user interface, the business logic, the
persistent data and the external location services. The
system unites station discovery, slot booking, live
availability and a rewards engine behind a single REST
API.

\subsection{Design Methodology and Core Objectives}

The engineering methodology focuses on four core
principles:

\begin{itemize}
  \item \textbf{Model Before Code:} The first weeks of the
    semester were spent on requirement analysis and UML
    modelling (use cases, activity, sequence and
    data-flow models), so that the frontend and backend
    could be built in parallel against a shared
    blueprint.
  \item \textbf{Loose Coupling:} The browser frontend
    communicates with the backend only through JSON REST
    calls and a WebSocket channel, so either side can
    change without breaking the other.
  \item \textbf{Real-Time by Default:} Port availability is
    pushed to the browser as it changes, instead of being
    refreshed manually.
  \item \textbf{Verifiable Quality:} Every use case has a
    written specification, and key workflows are covered
    by formal test cases with recorded results.
\end{itemize}

\subsection{Station Discovery and Smart Filtering Design}

A driver searches by location. The backend geocodes the
location and fetches nearby places through the Google
Maps Platform (Maps, Places and Directions APIs), merges
the result with station and charger records held in
MongoDB Atlas, and returns the stations with their price,
supported connectors and current availability. Smart
filters narrow the list by connector type, distance,
availability and price.

\subsection{Slot Booking and Real-Time Availability Design}

A booking captures the station, date, time slot,
connector type and vehicle number. The backend checks that
the slot is free, stores the booking record with a unique
booking identifier, and updates the station's free-port
count. The new count is broadcast to open browsers over a
WebSocket event so that availability badges change colour
without a page reload.

\subsection{EcoCoins Loyalty and Payment Design}

ChargEV rewards green mobility through an EcoCoins wallet.
At checkout the driver sees the order summary (station,
time window, charging fee and tax) and may redeem EcoCoins
against the payable amount; one EcoCoin is worth
INR~1.00. The redemption is validated against the wallet
balance, the discount is applied to the net total, and the
wallet is debited when payment is confirmed. After a
successful booking the driver is credited with a reward in
EcoCoins. Every debit and credit is written to a loyalty
ledger that is shown on the dashboard.

\subsection{Authentication and Session Design}

Drivers sign in with an e-mail address and password. The
backend verifies the credentials against the users
collection in MongoDB Atlas and, on success, returns a
JSON Web Token (JWT) together with the user profile. The
frontend keeps the token, the profile and a login flag in
the browser's \texttt{localStorage} and redirects to the
home page; the navigation bar then shows the user's
profile badge and a logout option. Invalid credentials
return an HTTP~401 response and a clear error message.

\subsection{System Decomposition and Multi-Tier Architecture}

The platform is organised into four decoupled tiers
(Table~\ref{tab:stack}):

\begin{enumerate}
  \item \textbf{Presentation Tier:} Responsive multi-page
    frontend built with HTML5, CSS3, Bootstrap~5 and
    JavaScript.
  \item \textbf{Application Tier:} Node.js and Express.js
    REST API that implements the authentication, station,
    booking, trip and rewards modules.
  \item \textbf{Persistence Tier:} MongoDB Atlas accessed
    through Mongoose schemas.
  \item \textbf{External Services Tier:} Google Maps,
    Places and Directions APIs for geocoding, place data
    and routing.
\end{enumerate}

\begin{table}[H]
  \centering
  \begin{tabularx}{\textwidth}{lX}
    \toprule
    \textbf{Layer} & \textbf{Technology} \\
    \midrule
    Frontend & HTML5, CSS3, Bootstrap 5, JavaScript (fetch API) \\
    Backend & Node.js, Express.js (REST API) \\
    Authentication & JSON Web Token (JWT), session kept in \texttt{localStorage} \\
    Database & MongoDB Atlas, accessed through Mongoose schemas \\
    Location intelligence & Google Maps, Places and Directions APIs \\
    Real-time updates & WebSocket events for live port-count updates \\
    Documentation & \LaTeX\ report, draw.io diagrams, MkDocs site \\
    \bottomrule
  \end{tabularx}
  \caption{Technology stack of ChargEV.}
  \label{tab:stack}
\end{table}

\section{System Modeling and Project Diagrams}

To provide a comprehensive blueprint of the ChargEV
system, several architectural and behavioural models were
developed across the semester. All diagrams were
cross-checked against the project proposal and against one
another.

\subsection{Project Scheduling and Timeline (Gantt Chart)}

Project implementation followed a structured 20-week
schedule covering requirement analysis, UI/UX design,
backend development, integration, testing and the final
report. Figure~\ref{fig:gantt} illustrates the project
timeline and milestones.

\diagram{gantt_chart.drawio.pdf}{Project Scheduling and Timeline (Gantt Chart).}{fig:gantt}
\clearpage

\subsection{UML Use Case Diagram}

The use case diagram (Figure~\ref{fig:usecase}) defines the
system boundary and its actors: the EV Driver (primary
actor), the Station Operator (secondary actor), and the
Location Services API and MongoDB Atlas (supporting
actors). It captures 17 use cases: searching nearby
stations, viewing station details, applying smart filters,
checking real-time availability, booking a charging slot,
navigating to a station, planning a multi-stop trip,
tracking a charging session, geocoding a location,
fetching places data, computing a route, storing a booking
record, updating station status, managing station
inventory, viewing the utilization dashboard, earning
EcoCoins, and viewing eco-impact and badges.

\diagram{use_case_diagram.drawio.pdf}{UML Use Case Diagram identifying system boundaries and actors.}{fig:usecase}

\subsection{Use Case Specifications}

The following pages give the specification template of
each use case: title, identifier, actors, description,
pre-conditions, task sequence, post-conditions,
modification history and authors.

\clearpage
\ucpage{UC_01_Search_Nearby_Stations.pdf}{Search Nearby Stations}
\ucpage{UC_02_View_Station_Details.pdf}{View Station Details}
\ucpage{UC_03_Apply_Smart_Filters.pdf}{Apply Smart Filters}
\ucpage{UC_04_Check_RealTime_Availability.pdf}{Check Real-Time Availability}
\ucpage{UC_05_Book_Charging_Slot.pdf}{Book Charging Slot}
\ucpage{UC_06_Navigate_to_Station.pdf}{Navigate to Station}
\ucpage{UC_07_Plan_MultiStop_Trip.pdf}{Plan Multi-Stop Trip}
\ucpage{UC_08_Track_Charging_Session.pdf}{Track Charging Session}
\ucpage{UC_09_Geocode_Location.pdf}{Geocode Location}
\ucpage{UC_10_Fetch_Places_Data.pdf}{Fetch Places Data}
\ucpage{UC_11_Compute_Route.pdf}{Compute Route}
\ucpage{UC_12_Store_Booking_Record.pdf}{Store Booking Record}
\ucpage{UC_13_Update_Station_Status.pdf}{Update Station Status}
\ucpage{UC_14_Manage_Station_Inventory.pdf}{Manage Station Inventory}
\ucpage{UC_15_View_Utilization_Dashboard.pdf}{View Utilization Dashboard}
\ucpage{UC_16_Earn_EcoCoins.pdf}{Earn EcoCoins}
\ucpage{UC_17_View_EcoImpact_and_Badges.pdf}{View Eco-Impact and Badges}

\subsection{UML Activity Diagram}

The activity diagram (Figure~\ref{fig:activity}) maps the
step-by-step workflow of a station search: the driver
enters a location and filters, the backend validates the
request, the location is resolved through the Maps API,
nearby stations are fetched through the Places API, slot
status is read from the database, the smart filters are
applied and the results are displayed on the map.

\diagram[0.85\textwidth]{activity_diagram.drawio.pdf}{UML Activity Diagram for the station search workflow.}{fig:activity}

\subsection{UML Activity Diagram (Swimlanes)}

The swimlane activity diagram (Figure~\ref{fig:swimlanes})
distributes the complete search--book--charge flow across
four responsibilities: the EV Driver/Frontend, the Backend
API (Node.js/Express), Location Services
(Maps/Places/Directions) and MongoDB Atlas.

\diagram{activity_diagram_swimlanes.drawio.pdf}{UML Activity Diagram with swimlanes.}{fig:swimlanes}

\subsection{UML Sequence Diagram}

The sequence diagram (Figure~\ref{fig:sequence}) shows the
chronological, message-level interaction between the User,
the Frontend web app, the Backend, the Database (MongoDB
Atlas) and the Google APIs while a driver searches for
stations and books a slot: from \texttt{GET /api/stations}
and the geocoding call, through the availability query, to
\texttt{POST /api/booking}, the booking confirmation and
the EcoCoins shown to the user.

\diagram{Sequence_Diagram.drawio.pdf}{UML Sequence Diagram: search and book a charging slot.}{fig:sequence}

\subsection{UML Collaboration Diagram}

The collaboration diagram (Figure~\ref{fig:collaboration})
shows the same interaction as a structural map of the
communication links between the User, Frontend, Backend,
Database and Google APIs, with each link numbered to show
the order of messages.

\diagram[0.9\textwidth]{Collaboration_Diagram.drawio.pdf}{UML Collaboration Diagram.}{fig:collaboration}

\subsection{UML Class Diagram}

The class diagram (Figure~\ref{fig:class}) defines the
domain model: \texttt{User}, \texttt{Operator},
\texttt{Station}, \texttt{Charger} and \texttt{Booking},
with their attributes, operations and multiplicities --
a user makes many bookings, a booking reserves one
charger, a station has many chargers, and an operator
manages many stations.

\diagram{class_diagram.pdf}{UML Class Diagram.}{fig:class}

\subsection{UML Component Diagram}

The component diagram (Figure~\ref{fig:component}) groups
the system into the browser frontend (station locator, trip
planner, login/signup, booking page and EcoCoins page), the
Node.js/Express backend (station, trip, auth, booking and
rewards modules), the Google APIs (Maps, Places and
Directions) and the MongoDB Atlas database, linked by HTTP
interfaces.

\diagram{component_diagram.pdf}{UML Component Diagram.}{fig:component}

\subsection{Data Flow Diagram (Level 0 - Context Level)}

The Level~0 DFD (Figure~\ref{fig:dfd0}) treats ChargEV as a
single process and shows its exchanges with three external
entities: the EV Driver (registration and login details,
search criteria, trip and booking requests in; stations,
routes, confirmations and EcoCoins out), the Station
Operator (station details and status updates in;
utilization reports out) and the Google Maps Platform
(geocode, place and route requests out; coordinates, place
data and routes in).

\diagram{dfd_level_0.drawio.pdf}{Level 0 Data Flow Diagram (Context Level).}{fig:dfd0}

\subsection{Data Flow Diagram (Level 1 - Functional Decomposition)}

The Level~1 DFD (Figure~\ref{fig:dfd1}) decomposes the
system into six processes -- Authenticate and Manage
Profile (1.0), Search and Filter Stations (2.0), Plan Trip
and Navigate (3.0), Book Slot and Track Session (4.0),
Manage Rewards and Eco-Impact (5.0) and Manage Stations and
Utilization (6.0) -- and four data stores: Users (D1),
Stations and Chargers (D2), Bookings (D3) and Rewards (D4).

\diagram{dfd_level_1.drawio.pdf}{Level 1 Data Flow Diagram illustrating sub-processes, data stores and internal information flow.}{fig:dfd1}

\subsection{UML Deployment Diagram}

The deployment diagram (Figure~\ref{fig:deployment}) maps
the software artifacts onto physical nodes: the client
device running the web browser, the application server
running Node.js and Express, the MongoDB Atlas cloud
database and the Google Maps Platform, together with the
communication protocols between them (HTTP/HTTPS, TLS and
HTTPS REST calls).

\diagram{deployment_diagram.drawio.pdf}{UML Deployment Diagram detailing nodes, execution environments and communication protocols.}{fig:deployment}

% =====================================================================
% CHAPTER 3: CODING AND IMPLEMENTATION
% =====================================================================

\chapter{Coding and Implementation}

This chapter details the software implementation of
ChargEV across its core subsystems: the responsive
frontend, authentication, station discovery, slot booking
with real-time availability, the EcoCoins loyalty and
payment flow, trip planning, and the backend API.

\section{Software Architecture and Subsystem Overview}

The code base is organised into a static, multi-page
frontend and a Node.js/Express backend. The frontend talks
to the backend through JSON REST calls made with the
browser \texttt{fetch} API, and receives live
port-availability updates over a WebSocket connection.
Table~\ref{tab:endpoints} lists the REST endpoints used by
the frontend.

\begin{table}[H]
  \centering
  \begin{tabularx}{\textwidth}{llX}
    \toprule
    \textbf{Method} & \textbf{Endpoint} & \textbf{Purpose} \\
    \midrule
    POST & \texttt{/api/auth/login} & Authenticate a user and return a JWT and profile \\
    GET  & \texttt{/api/stations}   & Return the list of charging stations \\
    POST & \texttt{/api/bookings}   & Create a slot booking for a station \\
    \bottomrule
  \end{tabularx}
  \caption{REST endpoints consumed by the frontend.}
  \label{tab:endpoints}
\end{table}

\todo{add any further endpoints implemented in server.js (register, my bookings, EcoCoins redeem/credit, trip planning) to Table~\ref{tab:endpoints}}

\shot[0.55]{code_project_structure.png}{Project folder structure of the ChargEV code base.}{fig:code-structure}

\section{User Interface and Frontend Implementation}

The frontend is a responsive multi-page application built
with HTML5, CSS3 and Bootstrap~5, with a dark, EV-themed
visual identity:

\begin{itemize}
  \item \textbf{Discovery Pages:} \texttt{index.html}
    (hero search, quick filters, nearest-station preview)
    and \texttt{station.html} (station directory with map
    and list view and live availability badges).
  \item \textbf{Booking Flow:} \texttt{booking.html} (station,
    date, time slot, connector and vehicle number) followed
    by \texttt{payment.html} (order summary, EcoCoins
    redemption and payment method).
  \item \textbf{Account Pages:} \texttt{login.html},
    \texttt{register.html} and \texttt{profile.html} (vehicle
    model and preferred connector).
  \item \textbf{Dashboard and History:}
    \texttt{dashboard.html} (charging session telemetry,
    saved bookings and EcoCoins wallet) and
    \texttt{mybookings.html} (a driver's scheduled sessions).
  \item \textbf{Information Pages:} \texttt{about.html} and
    \texttt{contact.html}.
\end{itemize}

Client-side logic is split by concern: \texttt{script.js}
handles the home page and the booking form,
\texttt{station.js} renders the station list,
\texttt{login.js} performs authentication and
\texttt{Contact.js} handles the contact form. Bootstrap's
grid keeps every screen responsive, and form controls
validate input through HTML5 attributes and JavaScript
event listeners.

\shot{ui_home.png}{Home Page: hero search, quick filters and nearest-station preview.}{fig:ui-home}
\shot{ui_stations.png}{Stations Page: station list with live port-availability badges.}{fig:ui-stations}
\shot{ui_booking.png}{Booking Page: slot, connector and vehicle selection.}{fig:ui-booking}
\shot{ui_payment.png}{Payment Page: order summary, EcoCoins redemption and payment method.}{fig:ui-payment}
\shot{ui_mybookings.png}{My Bookings: scheduled charging sessions.}{fig:ui-mybookings}
\shot{ui_dashboard.png}{Dashboard: charging session telemetry, saved bookings and EcoCoins wallet.}{fig:ui-dashboard}
\shot[0.6]{ui_login.png}{Login Page.}{fig:ui-login}
\shot[0.6]{ui_register.png}{Create Account Page.}{fig:ui-register}
\shot[0.6]{ui_profile.png}{Profile Page: vehicle model and connector preference.}{fig:ui-profile}

\section{Module 1: Authentication and Session Management}

Implemented in \texttt{login.html}, \texttt{login.js} and
the \texttt{User} model, this module verifies credentials
and keeps the driver's session. A user record holds the
name, e-mail (unique), password, phone, city, EV model and
preferred connector type (default \texttt{CCS2}), with
automatic timestamps (Listing~\ref{lst:user}).

\begin{lstlisting}[caption={Mongoose schema of the User model.},label={lst:user}]
const userSchema = new mongoose.Schema({
    name: { type: String, required: true },
    email: { type: String, required: true, unique: true },
    password: { type: String, required: true },
    phone: { type: String, default: '' },
    city: { type: String, default: '' },
    evModel: { type: String, default: '' },
    connectorType: { type: String, default: 'CCS2' }
}, { timestamps: true });
\end{lstlisting}

On submit, the frontend posts the credentials to
\texttt{/api/auth/login}. On success the returned token and
user object are stored in \texttt{localStorage} and the
user is redirected to the home page; on failure the error
returned by the server is shown in the login alert
(Listing~\ref{lst:login}).

\begin{lstlisting}[caption={Login request handling in \texttt{login.js}.},label={lst:login}]
const res = await fetch('http://localhost:5000/api/auth/login', {
  method: 'POST',
  headers: { 'Content-Type': 'application/json' },
  body: JSON.stringify({ email, password })
});
const data = await res.json();
if (res.ok) {
  localStorage.setItem('token', data.token);
  localStorage.setItem('user', JSON.stringify(data.user));
  localStorage.setItem('isLoggedIn', 'true');
  window.location.href = 'index.html';
}
\end{lstlisting}

\todo{describe how the server hashes passwords and signs the JWT, and add a code screenshot of the auth route}

\shot[0.8]{code_auth_route.png}{Backend authentication route (code screenshot).}{fig:code-auth}

\section{Module 2: Station Locator and Smart Filters}

The station list is fetched from \texttt{/api/stations} and
rendered as Bootstrap cards showing the station name, price
per hour, supported connector types and a \emph{Book Slot}
button that carries the station identifier to the booking
page (Listing~\ref{lst:stations}). If the list is empty or
the server cannot be reached, an informative message is
shown instead of failing silently.

\begin{lstlisting}[caption={Fetching and rendering stations in \texttt{station.js}.},label={lst:stations}]
const res = await fetch('http://localhost:5000/api/stations');
const stations = await res.json();
stationList.innerHTML = stations.map(station => `
  <div class="card mb-3 shadow-sm border-0 rounded-3">
    <div class="card-body">
      <h5 class="card-title text-success fw-bold">${station.name}</h5>
      <p class="card-text mb-1"><strong>Price:</strong> Rs.${station.pricePerHour}/hr</p>
      <a href="booking.html?stationId=${station._id}" class="btn btn-success btn-sm">Book Slot</a>
    </div>
  </div>`).join('');
\end{lstlisting}

\todo{describe the filters (distance, connector type, availability, price) and the Google Maps / Places integration as implemented, and add a code screenshot}

\shot[0.8]{code_stations_route.png}{Backend station route and filtering logic (code screenshot).}{fig:code-stations}

\section{Module 3: Slot Booking and Real-Time Availability}

A booking captures the station, date, time slot, connector
type and vehicle number, and is submitted to
\texttt{/api/bookings} (Listing~\ref{lst:booking}). On
success the user is taken to the bookings page; on failure
the server's message is displayed. Each booking is stored
using the \texttt{Booking} schema
(Listing~\ref{lst:bookingmodel}), which holds a unique
\texttt{bookingId} and records the connector, duration,
time slot and total price.

\begin{lstlisting}[caption={Booking submission in \texttt{script.js}.},label={lst:booking}]
const bookingData = {
    stationId: selectedId,
    stationName: activeStation ? activeStation.name : 'EV Hub',
    date: bookingDate.value,
    timeSlot: timeSlot.value,
    connectorType: document.getElementById('connectorType').value,
    vehicleNumber: document.getElementById('vehicleNumber').value.trim()
};
const response = await fetch('/api/bookings', {
    method: 'POST',
    headers: { 'Content-Type': 'application/json' },
    body: JSON.stringify(bookingData)
});
\end{lstlisting}

\begin{lstlisting}[caption={Mongoose schema of the Booking model.},label={lst:bookingmodel}]
const BookingSchema = new mongoose.Schema({
    bookingId: { type: String, required: true, unique: true },
    stationId: { type: String, required: true },
    connector: { type: String, required: true },
    durationMinutes: { type: Number, required: true },
    timeSlot: { type: String, required: true },
    totalPrice: { type: String, required: true },
    createdAt: { type: Date, default: Date.now }
});
\end{lstlisting}

\textbf{Live availability.} The station cards subscribe to a
\texttt{port\_count\_updated} WebSocket event. When the
server reports a new number of free ports for a station,
the badge on the corresponding card is updated without a
page reload: green when more than one port is free, amber
when exactly one is free and red when none is free
(Listing~\ref{lst:ws}).

\begin{lstlisting}[caption={Live port-count updates in \texttt{script.js}.},label={lst:ws}]
socket.on('port_count_updated', (data) => {
    const badge = document.getElementById(`badge-${data.stationId}`);
    if (badge) {
        badge.innerText = `${data.availablePorts} ${data.availablePorts === 1 ? 'Port' : 'Ports'} Free`;
        badge.className = `... badge ${data.availablePorts > 1 ? 'bg-success'
            : (data.availablePorts === 1 ? 'bg-warning text-dark' : 'bg-danger')} ...`;
    }
});
\end{lstlisting}

\todo{add a code screenshot of the booking route (slot-conflict check, bookingId generation, port-count update)}

\shot[0.8]{code_booking_route.png}{Backend booking route (code screenshot).}{fig:code-booking}

\section{Module 4: EcoCoins Loyalty and Payment Checkout}

Implemented across \texttt{payment.html},
\texttt{booking.html} and \texttt{dashboard.html}, this
module manages the EcoCoins lifecycle during checkout:

\begin{itemize}
  \item \textbf{Order summary:} the payment page shows the
    station, time window, charging fee, tax and the initial
    total payable amount.
  \item \textbf{Redemption:} the driver enters the EcoCoins
    to redeem; the entry is validated against the wallet
    balance, converted to a discount at
    1~EcoCoin~=~INR~1.00, and the net total is updated.
  \item \textbf{Payment and debit:} after choosing a
    payment method (UPI/QR code or card), confirming
    payment debits the redeemed EcoCoins and creates the
    confirmed booking record.
  \item \textbf{Reward:} on success the driver is shown a
    ``Payment Successful'' screen and credited with a
    green-mobility reward in EcoCoins.
  \item \textbf{Ledger:} the dashboard wallet widget shows
    the updated balance and logs each redemption and
    reward in the loyalty ledger.
\end{itemize}

This behaviour is verified by test case TC\_CHARGEV\_01
(Section~\ref{sec:tc01}).

\todo{add a code screenshot of the EcoCoins redeem / credit logic and the wallet ledger model}

\shot[0.8]{code_ecocoins.png}{EcoCoins redemption, reward and ledger logic (code screenshot).}{fig:code-eco}

\section{Module 5: Trip Planning and Navigation}

\todo{describe the trip planning module as implemented (Google Directions API, charging stops along the route); if it was not implemented, state this here and move it to Limitations}

\shot[0.8]{code_trip_planner.png}{Trip planning module (code screenshot).}{fig:code-trip}

\section{Module 6: Station Operator Tools}

\todo{describe the operator functions as implemented (add station, update charger status, utilization dashboard); if not implemented, state this and move it to Limitations}

\section{Module 7: REST API Controller and Database Integration}

The Express server exposes the REST endpoints of
Table~\ref{tab:endpoints}, connects to MongoDB Atlas through
Mongoose and broadcasts port-count changes over WebSocket.
During development the frontend expects the server on
\texttt{localhost:5000}.

\todo{add a code screenshot of server.js (Express setup, CORS, MongoDB connection) and the .env variable names (without secret values)}

\shot[0.8]{code_server.png}{Express server setup and MongoDB Atlas connection (code screenshot).}{fig:code-server}

% =====================================================================
% CHAPTER 4: RESULTS AND EVALUATION
% =====================================================================

\chapter{Results and Evaluation}

\section{Testing Methodology}

Testing was conducted across three tiers:

\begin{itemize}
  \item \textbf{Design Validation:} The UML and DFD models
    were cross-checked against the proposal and against
    each other (for example, each Level~1 process maps to
    one or more use cases, and each data store maps to a
    class or collection of the backend).
  \item \textbf{Functional (Black-Box) Testing:} Each
    formal test case exercises one use case through the
    user interface and compares the observed behaviour
    with the expected system response, recording a
    Pass/Fail status for every step.
  \item \textbf{Integration Testing:} The frontend was
    exercised against the running Express backend and
    the MongoDB Atlas database, including the
    authentication service and the loyalty ledger.
\end{itemize}

\section{Software Engineering Test Cases (SE Lab Deliverable)}

To fulfil the Software Engineering verification
requirements, the formalized test cases of the ChargEV test
specification are presented below. Each test case is
structured on its own dedicated page.

\clearpage
\subsection{Test Case (Doc:TC\_CHARGEV\_01)}
\label{sec:tc01}

\begin{center}
\footnotesize
\renewcommand{\arraystretch}{1.05}
\begin{tabularx}{\textwidth}{|p{2.8cm}|X|p{2.8cm}|X|}
\hline
\textbf{Test Case \#:} & TC\_CHARGEV\_01 & \textbf{Test Case Name:} & Process EcoCoins Debit (Redemption Discount) and Credit (Reward Accrual) on Booking \\
\hline
\textbf{System:} & ChargEV -- Smart EV Charging Station Platform & \textbf{Subsystem:} & Loyalty \& Payment Gateway Subsystem (payment.html / booking.html / UC-16) \\
\hline
\textbf{Designed by:} & Dhruv and Srishti & \textbf{Design Date:} & 20/11/2026 \\
\hline
\textbf{Executed by:} & Pehal and Briti Arora & \textbf{Execution Date:} & 20/11/2026 \\
\hline
\textbf{Page:} & 1 of 1 & & \\
\hline
\multicolumn{4}{|p{0.96\textwidth}|}{\textbf{Short Description:} \newline Verify the transactional lifecycle of EcoCoins during slot checkout: validating wallet balance, debiting redeemed EcoCoins to apply a tariff discount, completing payment, and crediting newly earned green mobility reward EcoCoins upon booking confirmation.} \\
\hline
\end{tabularx}
\end{center}

\vspace{-0.4em}
\noindent\textbf{Pre-conditions:}
\begin{enumerate}\footnotesize
    \setlength{\itemsep}{1pt}
    \setlength{\parskip}{0pt}
    \item User is authenticated (e.g.\ rohini@example.com) and has navigated to the checkout page \texttt{payment.html}.
    \item User's EcoCoins wallet balance is pre-loaded with 120 EcoCoins (conversion: 1 EcoCoin = INR 1.00).
    \item A charging slot is selected with an initial total payable amount of INR 354.00 (tariff INR 300.00 + tax INR 54.00).
\end{enumerate}

\vspace{-0.4em}
\noindent\textbf{Execution Steps:}
\begin{center}
\footnotesize
\renewcommand{\arraystretch}{1.02}
\setlength{\tabcolsep}{3pt}\begin{tabularx}{\textwidth}{|>{\centering\arraybackslash}p{0.9cm}|>{\raggedright\arraybackslash}p{3.3cm}|>{\raggedright\arraybackslash}X|>{\centering\arraybackslash}p{1.1cm}|>{\centering\arraybackslash\scriptsize}p{2.3cm}|}
\hline
\textbf{Step} & \textbf{Action} & \textbf{Expected System Response} & \textbf{Status} & \textbf{Comment} \\
\hline
1 & Review the order summary on the checkout page (\texttt{payment.html}). & Station details, time window, charging fee (INR 300.00), tax (INR 54.00) and initial total amount of INR 354.00 are displayed. & \textbf{PASS} & {\scriptsize Checkout summary verified.} \\
\hline
2 & Locate the `Redeem EcoCoins' field and enter 50 EcoCoins to redeem for this booking. & System validates the input against the wallet balance (120 $\geq$ 50), calculates a discount of INR 50.00 and updates the net total amount to INR 304.00. & \textbf{PASS} & {\scriptsize Discount calculation confirmed.} \\
\hline
3 & Select a payment method (UPI / QR code or card) and click `Pay INR 304.00 \& Confirm'. & System processes the transaction, debits 50 EcoCoins from the wallet (remaining balance: 70 EcoCoins) and creates a confirmed booking record. & \textbf{PASS} & {\scriptsize Wallet debit transaction verified.} \\
\hline
4 & Inspect the booking confirmation screen and reward summary. & System displays `Payment Successful!', awards +15 EcoCoins green mobility cashback reward and updates the final wallet balance to 85 EcoCoins. & \textbf{PASS} & {\scriptsize Reward crediting confirmed.} \\
\hline
5 & Navigate to the Dashboard (\texttt{dashboard.html}) and inspect the EcoCoins wallet widget. & Dashboard reflects the updated balance of 85 EcoCoins and logs the ledger records `-50 EcoCoins (Redeemed)' and `+15 EcoCoins (Reward)'. & \textbf{PASS} & {\scriptsize Ledger persistence verified.} \\
\hline
\end{tabularx}
\end{center}

\vspace{-0.4em}
\noindent\textbf{Post-conditions:}
\begin{enumerate}\footnotesize
    \setlength{\itemsep}{1pt}
    \setlength{\parskip}{0pt}
    \item User wallet balance is accurately updated from 120 to 85 EcoCoins (net change: $-35$ EcoCoins).
    \item The booking is recorded as paid with the discounted amount (INR 304.00) in the database.
    \item EcoCoins debit and credit transactions are permanently logged in the user's loyalty ledger.
\end{enumerate}

\clearpage
\subsection{Test Case (Doc:TC\_CHARGEV\_02)}
\label{sec:tc02}

\begin{center}
\footnotesize
\renewcommand{\arraystretch}{1.05}
\begin{tabularx}{\textwidth}{|p{2.8cm}|X|p{2.8cm}|X|}
\hline
\textbf{Test Case \#:} & TC\_CHARGEV\_02 & \textbf{Test Case Name:} & Authenticate EV Driver with Credential Verification and Session Management \\
\hline
\textbf{System:} & ChargEV -- Smart EV Charging Station Platform & \textbf{Subsystem:} & User Authentication \& Session Management (login.html / login.js / user.js) \\
\hline
\textbf{Designed by:} & Dhruv and Srishti & \textbf{Design Date:} & 20/11/2026 \\
\hline
\textbf{Executed by:} & Pehal and Briti Arora & \textbf{Execution Date:} & 20/11/2026 \\
\hline
\textbf{Page:} & 1 of 1 & & \\
\hline
\multicolumn{4}{|p{0.96\textwidth}|}{\textbf{Short Description:} \newline Verify the user authentication workflow: invalid credential rejection with error alert display, input field validation, successful login with valid credentials, session token persistence in localStorage, and redirection to the application home page.} \\
\hline
\end{tabularx}
\end{center}

\vspace{-0.4em}
\noindent\textbf{Pre-conditions:}
\begin{enumerate}\footnotesize
    \setlength{\itemsep}{1pt}
    \setlength{\parskip}{0pt}
    \item The ChargEV web application is accessible at \texttt{login.html}.
    \item The backend authentication service (\texttt{POST /api/auth/login}) and the MongoDB Atlas users collection are active.
    \item A verified test user exists in the database (e.g.\ e-mail: rohini@example.com, password: Password@123).
    \item Browser \texttt{localStorage} is cleared of prior authentication tokens (\texttt{token}, \texttt{isLoggedIn}).
\end{enumerate}

\vspace{-0.4em}
\noindent\textbf{Execution Steps:}
\begin{center}
\footnotesize
\renewcommand{\arraystretch}{1.02}
\setlength{\tabcolsep}{3pt}\begin{tabularx}{\textwidth}{|>{\centering\arraybackslash}p{0.9cm}|>{\raggedright\arraybackslash}p{3.3cm}|>{\raggedright\arraybackslash}X|>{\centering\arraybackslash}p{1.1cm}|>{\centering\arraybackslash\scriptsize}p{2.3cm}|}
\hline
\textbf{Step} & \textbf{Action} & \textbf{Expected System Response} & \textbf{Status} & \textbf{Comment} \\
\hline
1 & Navigate to the login interface (\texttt{login.html}). & Login card is displayed with e-mail field, password field, `Remember me' checkbox and `Sign In' button; the error alert \texttt{\#login-error} is hidden. & \textbf{PASS} & {\scriptsize Initial UI render verified.} \\
\hline
2 & Enter the registered e-mail `rohini@example.com' and the incorrect password `WrongPass!99', then click `Sign In'. & System sends \texttt{POST /api/auth/login}, receives HTTP 401 Unauthorized, shows the message `Invalid Credentials'; \texttt{localStorage} remains empty. & \textbf{PASS} & {\scriptsize Negative test / invalid credential handling verified.} \\
\hline
3 & Clear the password field, leave it blank and click `Sign In'. & HTML5 client-side validation triggers on \texttt{\#login-password}; the browser halts submission and prompts `Please fill out this field'. & \textbf{PASS} & {\scriptsize Client form validation confirmed.} \\
\hline
4 & Enter the correct password `Password@123' and click `Sign In'. & System sends \texttt{POST /api/auth/login} and receives HTTP 200 OK with a JWT token and user profile object in the JSON response. & \textbf{PASS} & {\scriptsize Backend authentication confirmed.} \\
\hline
5 & Inspect browser \texttt{localStorage} and observe the redirection. & System stores \texttt{token}, \texttt{user} and \texttt{isLoggedIn='true'}; \texttt{\#login-error} remains hidden; the browser redirects to \texttt{index.html}. & \textbf{PASS} & {\scriptsize Session persistence confirmed.} \\
\hline
6 & Inspect the navigation bar on \texttt{index.html}. & Header reflects the authenticated state (the `Login/Register' buttons are replaced by a user profile badge / logout option). & \textbf{PASS} & {\scriptsize Authenticated navbar state verified.} \\
\hline
\end{tabularx}
\end{center}

\vspace{-0.4em}
\noindent\textbf{Post-conditions:}
\begin{enumerate}\footnotesize
    \setlength{\itemsep}{1pt}
    \setlength{\parskip}{0pt}
    \item Driver session is established with an authenticated JWT token stored in client storage.
    \item Failed login attempts are intercepted with clear UI error feedback and do not grant unauthorized access.
\end{enumerate}

\clearpage
\section{Test Execution Summary}

Table~\ref{tab:testsummary} summarises the outcome of the
executed test cases.

\begin{table}[H]
\centering
\caption{Test Execution Summary}
\label{tab:testsummary}
\begin{tabular}{llccc}
\toprule
\textbf{Test Case} & \textbf{Use case / subsystem} & \textbf{Steps} & \textbf{Passed} & \textbf{Failed} \\
\midrule
TC\_CHARGEV\_01 & EcoCoins checkout (UC-16) & 5 & 5 & 0 \\
TC\_CHARGEV\_02 & Authentication \& session & 6 & 6 & 0 \\
\midrule
\textbf{Total} & & \textbf{11} & \textbf{11} & \textbf{0} \\
\bottomrule
\end{tabular}
\end{table}

\todo{if more test cases are added (search, booking, availability), add their rows here}

\section{Performance Metrics}

The proposal defines the metrics used to judge ChargEV's
booking workflow (Table~\ref{tab:metrics}). They provide the
framework against which the system is measured as it is
used with live station data.

\begin{table}[H]
\centering
\caption{Evaluation metrics framework for the ChargEV booking workflow.}
\label{tab:metrics}
\begin{tabularx}{\textwidth}{lX}
  \toprule
  \textbf{Metric} & \textbf{Definition / Target} \\
  \midrule
  Time-to-Charge (TTC) & Time from opening the app to
    starting a charge at a booked slot; the target is a
    lower median TTC than a no-booking, walk-in
    baseline. \\
  Booking Conversion Rate & Percentage of station
    searches that convert into a confirmed booking. \\
  Availability Data Accuracy & Agreement between the
    displayed station status and the actual status
    observed on arrival. \\
  User Engagement & Repeat-usage rate and EcoCoins
    accumulated per active user. \\
  \bottomrule
\end{tabularx}
\end{table}

\todo{if any measurements were taken (for example API response time or page load time), add a Target vs.\ Achieved table here}

\section{Analysis}

Modelling before coding gave the frontend and backend a
shared contract: every screen matched the data returned by
the backend, so integration needed little UI refactoring.
The Swimlane Activity diagram and the Level~1 DFD fixed the
responsibility boundaries between the frontend, the Express
API, MongoDB Atlas and the Google APIs before the backend
was written. The executed test cases confirm two
business-critical workflows end to end: secure sign-in with
clear error handling and session persistence (all 6 steps
passed), and the EcoCoins checkout in which redemption,
payment, reward and ledger records stay consistent (all 5
steps passed). The live availability badge removes the main
uncertainty a driver faces on arrival -- whether a port is
free -- and the booking form removes the wait by reserving
a slot in advance.

% =====================================================================
% CHAPTER 5: CONCLUSION
% =====================================================================

\chapter{Conclusion and Summary}

\section{Summary of Work}

This project designed, implemented and validated
\textbf{ChargEV}, a web platform for discovering, booking
and using EV charging stations. It fulfilled the objectives
set in Chapter~1:

\begin{enumerate}
  \item \textbf{Station Locator:} A responsive frontend and
    REST API that list stations with price, connectors
    and live availability.
  \item \textbf{Real-Time Availability and Booking:} Slot
    booking with unique booking identifiers and live
    port-count updates over WebSocket.
  \item \textbf{Location Intelligence:} Integration with
    the Google Maps Platform for geocoding, places and
    routing.
  \item \textbf{Engagement Layer:} An EcoCoins wallet with
    redemption at checkout, reward on booking and a
    persistent loyalty ledger.
  \item \textbf{Software Quality:} Seventeen use case
    specifications, a full UML and DFD model set and
    formally documented test cases, all of which passed.
\end{enumerate}

\todo{confirm that every item in this summary matches what is actually implemented (Google Maps integration, live port updates, trip planning) and edit accordingly}

\section{Key Findings}

The project yielded three principal findings:

\begin{enumerate}
  \item \textbf{Model-First Development Pays Off:} Use
    cases, sequence diagrams and data-flow models
    written before coding let frontend and backend
    progress in parallel and integrate cleanly.
  \item \textbf{Live Availability and Advance Booking
    Address the Core Pain Points:} Knowing whether a port
    is free, and reserving it ahead of time, removes the
    uncertainty and waiting that define the current EV
    charging experience.
  \item \textbf{A Rewards Ledger Must Stay Consistent with
    Payments:} Debiting redemptions, crediting rewards
    and logging both in one ledger keeps the wallet,
    booking and dashboard in agreement, as confirmed by
    TC\_CHARGEV\_01.
\end{enumerate}

\section{Software Engineering Contributions}

From a software engineering perspective, ChargEV exemplifies
the full life cycle:

\begin{itemize}
  \item \textbf{Strict Modularity and Decoupling:} A clear
    separation between the presentation tier, the Express
    API, the MongoDB persistence tier and the external
    Google services.
  \item \textbf{Comprehensive UML Specifications:} The
    system was formalised through seventeen detailed use
    case specifications, activity diagrams with swimlanes,
    sequence and collaboration diagrams, a domain class
    diagram, a component diagram, Level~0 and Level~1
    data-flow diagrams and a physical deployment model.
  \item \textbf{Verified Quality and Collaboration:} Formal
    test cases with recorded results, version-controlled
    collaboration on GitHub and a weekly development
    journal.
\end{itemize}

\section{Limitations}

Despite meeting its objectives, ChargEV has several
limitations:

\begin{enumerate}
  \item \textbf{Prototype Authentication Storage:} The JWT
    is kept in \texttt{localStorage}, which is acceptable
    for a prototype but should be hardened (for example
    HTTP-only cookies) for production.
  \item \textbf{No Charger Hardware Integration:}
    Availability depends on data entered by operators and
    booking events; there is no live connection to
    physical chargers.
  \item \textbf{Simulated Payments:} The checkout flow
    does not process real financial transactions.
  \item \textbf{Limited Test Coverage:} Formal test cases
    cover authentication and EcoCoins checkout; automated
    unit and integration test suites are not yet in
    place.
  \item \todo{list any modules from Chapter~3 that were not fully implemented, e.g. trip planning, operator tools}
\end{enumerate}

\section{Future Work and Recommendations}

To build upon this foundation, future work will focus on
three key directions:

\begin{enumerate}
  \item \textbf{Real Payment Gateway and Charger
    Integration:} Connecting a live payment provider and
    charger status through the OCPP protocol.
  \item \textbf{Smarter Trip Planning:} Battery-range and
    charging-speed aware route optimisation with
    automatic stop selection.
  \item \textbf{Mobile Clients and Notifications:} Native
    apps with booking reminders, together with automated
    tests and a CI/CD pipeline.
\end{enumerate}

\section{Final Remarks}

ChargEV demonstrates that discovery, availability, booking
and engagement can be brought together in a single
dependable platform for EV drivers. By uniting live
availability, advance reservation, an EcoCoins rewards
loop and rigorous software engineering practice, this
project delivers a solid base for taking EV charging from a
guessing game to a predictable, rewarding experience.

% ===== BIBLIOGRAPHY =====
\nocite{*}
\printbibliography[heading=bibintoc]

\end{document}
